\section{Object content}

The object model is organised around a three-level domain hierarchy:
\emph{POI types} are aggregated into \emph{services}, and services are
aggregated into \emph{capabilities}. All spatial results are evaluated for an
origin transport-network node located within the study area. A POI may have
more than one network access point when its geometry is linear or polygonal.

\subsection*{Input objects}
\begin{enumerate}
    \item \textbf{Boundary of extraction.}
    The boundary object defines the spatial extent associated with an analysis
    run. Its mandatory fields are a polygon or multipolygon geometry and its
    coordinate reference system. It also carries a stable boundary identifier,
    a human-readable area name, and any source attributes required to identify
    the represented administrative or analytical unit. The boundary geometry
    determines the area in which origin transport-network nodes are evaluated.
    The POI extraction extent is currently coincident with this geometry; a
    separate buffered extraction geometry is specified by the cut-off threshold
    object described below.

    \item \textbf{POI shapefile collection.}
    The POI input is a collection of point, line, and polygon feature sets that
    share one logical POI schema. Each POI record contains:
    \begin{itemize}
        \item a source identity composed of an element category and a source
        identifier;
        \item a geometry and coordinate reference system;
        \item a POI classification label used to associate the record with a
        configured POI type;
        \item a tag family identifying the principal classification attribute;
        \item optional descriptive attributes, including name, address,
        operator, opening hours, accessibility attributes, and source-specific
        thematic tags.
    \end{itemize}
    Point geometries represent one candidate location. Line and polygon
    geometries represent spatially extended POIs and therefore admit multiple
    candidate access locations along their vertices and boundaries. The POI
    classification label is matched against the labels declared by the POI type
    configuration.

    \item \textbf{GTFS feed.}
    The GTFS object is a relational public-transport timetable dataset. Its
    principal entities and relations are:
    \begin{itemize}
        \item \emph{agency}: operator identity, name, timezone, language, and
        contact information;
        \item \emph{stop}: stop identifier, name, longitude, latitude, location
        type, parent station, and accessibility attributes;
        \item \emph{route}: route identifier, agency reference, route names,
        transport type, and display attributes;
        \item \emph{trip}: trip identifier, route reference, service-calendar
        reference, direction, headsign, shape reference, and accessibility
        attributes;
        \item \emph{stop time}: trip reference, stop reference, stop sequence,
        arrival time, departure time, and boarding or alighting constraints;
        \item \emph{calendar} and \emph{calendar exception}: service identifier,
        active weekdays, validity interval, and date-specific additions or
        removals;
        \item \emph{shape}: shape identifier and an ordered sequence of
        geographic points describing a trip path.
    \end{itemize}
    Optional feed entities include transfers, pathways, fares, booking rules,
    and feed metadata. The trip, stop-time, calendar, and stop relations jointly
    define which public-transport connections are available at the configured
    departure date and time.
\end{enumerate}

\subsection*{Configuration objects}
\begin{enumerate}
    \item \textbf{General configuration.}
    The general configuration identifies one analysis run and binds its spatial,
    transport, temporal, and visualisation objects. Its fields are grouped as
    follows:
    \begin{itemize}
        \item \emph{study-area identity}: study-city key, city name, city slug,
        analysis-run identifier, boundary-selection flag, and boundary
        reference;
        \item \emph{POI source selection}: source type and references to the
        point, line, and polygon POI collections;
        \item \emph{transport scope}: the walking, cycling, driving, and public
        transport modes; the GTFS feed references; the public-transport
        departure date and time; and the coefficient controlling the relative
        contribution of in-vehicle time to public-transport impedance;
        \item \emph{evaluation scope}: optional limits on the set of origin
        nodes and POIs included in the analysis;
        \item \emph{visualisation specification}: selected capability field,
        classification count, colour ramp, point and grid layer settings, grid
        cell size, grid opacity, and maximum grid-cell count.
    \end{itemize}

    \item \textbf{POI types configuration.}
    Each POI type configuration record defines one conceptual class of
    destination. Its fields are:
    \begin{itemize}
        \item \texttt{poi\_type}: unique POI type identifier;
        \item \texttt{decay\_coefficient}: impedance value, in minutes, at
        which access to the POI type has a decay value of \(0.5\);
        \item \texttt{choquet\_interactions}: ordered interaction coefficients
        between this POI type and the other POI types belonging to the same
        service; the self-interaction position is null;
        \item \texttt{tags}: one conjunctive tag clause or a disjunction of tag
        clauses used to classify source POIs;
        \item \texttt{labels}: accepted source labels used to classify POIs
        supplied as shapefile records.
    \end{itemize}
    The interaction-vector order is defined by the ordered POI type list of the
    associated service.

    \item \textbf{Service configuration.}
    Each service configuration record defines an ordered aggregation of POI
    types. Its fields are:
    \begin{itemize}
        \item \texttt{service}: unique service identifier;
        \item \texttt{poi\_types}: ordered list of constituent POI type
        identifiers;
        \item \texttt{choquet\_capacity}: ordered singleton capacities
        associated with the constituent POI types;
        \item \texttt{contribution\_coefficient}: ordered coefficients
        controlling how repeated accessible instances of each POI type
        contribute to its accessibility value.
    \end{itemize}
    The three ordered arrays have equal cardinality. A service record is valid
    only when every referenced POI type exists in the POI type configuration.

    \item \textbf{Capability configuration.}
    Each capability configuration record defines an ordered aggregation of
    services. Its fields are:
    \begin{itemize}
        \item \texttt{capability}: unique capability identifier;
        \item \texttt{services}: ordered list of constituent service
        identifiers;
        \item \texttt{electre\_weight}: ordered service-weight vector aligned
        with the service list;
        \item \texttt{enabled}: Boolean indicating whether the capability is
        included in the result set.
    \end{itemize}
    The configured capability set comprises restorativeness, nutrition, and
    care. Every referenced service must exist in the service configuration.
\end{enumerate}

\subsection*{Intermediate objects}
\begin{enumerate}
    \item \textbf{Impedance matrix.}
    The impedance object relates origin transport-network nodes to candidate POI
    access locations. Its origin dimension contains the origin node identifier
    and coordinates. Its destination dimension contains a destination access
    identifier, snapped access coordinates, and the associated POI type. For
    each origin--destination pair, the object may contain walking, cycling,
    driving, and public-transport impedance values expressed in minutes or an
    explicit unreachable value.

    Walking, cycling, and driving values are aligned vectors: position \(i\) in
    each modal vector refers to the same candidate POI access location. Walking
    impedance is associated with the walkability score of the selected walking
    path. Public-transport impedance is represented as an origin-by-destination
    matrix and is accompanied by origin and destination index mappings. The
    impedance object also carries the departure timestamp and signatures of the
    ordered origin and destination sets so that the represented routing problem
    is unambiguous.

    \item \textbf{Cut-off threshold and POI extraction buffer.}
    This is a required derived object that is not yet represented independently
    in the current object set. It separates the analysis boundary from the
    larger area from which potentially relevant POIs and routing destinations
    are selected. Its fields shall be:
    \begin{itemize}
        \item a minimum relevant decay value \(d_{\min}\), with
        \(0 < d_{\min} < 1\);
        \item the POI-type-specific maximum relevant impedance
        \(I_{\max,p}=-\ln(d_{\min})c_p/\ln(2)\), where \(c_p\) is the configured
        decay coefficient of POI type \(p\);
        \item the mode-specific maximum relevant distance derived from each
        maximum impedance and the corresponding modal travel relation;
        \item the conservative extraction-buffer distance selected from the
        mode- and POI-type-specific limits;
        \item the source boundary geometry and the resulting buffered extraction
        geometry;
        \item the set of POI types and modes covered by the threshold.
    \end{itemize}
    The resulting buffered geometry is the spatial selection object for POIs,
    while origin nodes remain associated with the unbuffered analysis boundary.

    \item \textbf{Bus routing result.}
    A bus routing result represents the selected public-transport connection for
    one origin and one destination access location at a specified departure
    time. A route record contains the origin identifier, destination identifier,
    departure time, access time, waiting time, in-vehicle time, transfer time,
    egress time, route sequence, number of rides, and total time. The associated
    impedance record contains travel time excluding waiting, waiting time,
    weighted impedance, and the route sequence. The result set carries the
    departure timestamp and ordered-origin and ordered-destination signatures.

    \item \textbf{Walking, cycling, and driving routing results.}
    A non-public-transport routing result is defined for one origin node and one
    POI type. It contains the origin coordinates, an ordered list of candidate
    POI coordinates, and three equally aligned impedance vectors for walking,
    cycling, and driving. An entry is null when the corresponding candidate is
    unreachable by that mode. The object additionally contains the path-level
    walkability score aligned with each walking impedance. Results are grouped
    by service while retaining the POI type identifier of every entry.

    \item \textbf{Walkability scores for edges.}
    The edge walkability object is a relation from a directed network-edge
    identifier \((u,v,k)\) to a walkability score in the interval from \(1\) to
    \(5\). It contains a graph identity specifying the network snapshot to which
    the scores apply. A path walkability value is a length-weighted aggregation
    of the edge scores belonging to that path. Each edge score is conceptually
    associated with observations at the start, midpoint, and end of the edge.

    \item \textbf{Accessibility results.}
    An accessibility result is defined for one origin node. Its fields are the
    node identifier, latitude, longitude, and a mapping from each service to an
    ordered collection of POI-type accessibility records. Each POI-type record
    contains a POI type identifier and an accessibility value in \([0,1]\).

    For each candidate POI access location, the four modal impedance values are
    transformed into modal decay values and combined into one multimodal
    route-resource-availability value. The ordered collection of candidate
    values for a POI type is then reduced to the POI-type accessibility value.
    The result preserves service membership without yet aggregating the POI
    types into a service score.

    \item \textbf{Service results.}
    A service result is defined for one origin node. Its fields are the node
    identifier, latitude, longitude, and a mapping from every configured service
    identifier to one normalized service score in \([0,1]\). Each score
    aggregates the accessibility values of the ordered POI types declared by
    that service, using the associated singleton capacities and pairwise
    interactions.
\end{enumerate}

\subsection*{Output objects}
\begin{enumerate}
    \item \textcolor{blue}{\textbf{Capability results.}}
    A capability result is defined for one origin node and one enabled
    capability. Its fields are the node identifier, latitude, longitude, the
    normalized capability score in \([0,1]\), and the complete set of service
    scores that constitute that capability. The service-score fields retain the
    service identifiers, allowing each capability value to be decomposed into
    its direct service components.

    A run-level capability summary contains the study-city identifier,
    analysis-run identifier, and the mean score of each capability over all
    evaluated origin nodes.

    \item \textbf{Visualisation map.}
    The visualisation map is a spatial object composed of two related layers:
    \begin{itemize}
        \item the \emph{capability point layer}, containing one point per
        evaluated origin node, with node identifier, longitude, latitude, point
        geometry, and one field for each available capability result;
        \item the \emph{capability grid layer}, containing regular hexagonal
        cell geometry, column and row indices, the mean of each capability over
        the nodes contained by the cell, a selected display mean, a data-
        availability indicator, the selected capability field, and the cell's
        metric dimensions.
    \end{itemize}
    The map specification associates the spatial layers with the selected
    capability field, classification count, colour ramp, opacity, and optional
    basemap reference.

    \item \textcolor{blue}{provenance interface}
\end{enumerate}
